% Generated by paper/tools/render_experiment_figures.py; do not edit by hand.
\begin{table*}[t]
\centering
\small
\caption{Current A100 80GB matched baseline comparisons. Ratios are baseline median divided by cuButterfly median; values above one favor cuButterfly.}
\label{tab:current-results}
\begin{tabular}{llrrrr}
\hline
Class & Baseline & Operator & Cells & Geo. mean & Ours faster \\
\hline
external & Dao & FWHT & 8 & 2.27793 & 6 \\
external & GPU-NTT & NTT & 6 & 1.12917 & 4 \\
external & VkFFT & FFT & 18 & 0.99278 & 8 \\
external & cuFFT & FFT & 29 & 0.97893 & 11 \\
in-tree control & in-tree radix2 & FFT & 29 & 2.95929 & 29 \\
in-tree control & in-tree radix2 & FWHT & 12 & 2.18311 & 12 \\
in-tree control & in-tree radix2 & structured-2x2 & 6 & 1.61207 & 6 \\
in-tree control & in-tree radix2 & subset-zeta & 5 & 1.43396 & 5 \\
in-tree control & in-tree radix2 & superset-zeta & 5 & 1.43305 & 5 \\
in-tree control & in-tree radix2 & xor-zeta & 1 & 1.39377 & 1 \\
\hline
\end{tabular}
\parbox{0.98\linewidth}{\footnotesize A100 80GB PCIe, SM80; 100 warmups, 100 repetitions, three process trials. Each cell contributes the median of its three kernel-ms samples before the geometric mean is taken. The in-tree radix2 rows are ordinary-core controls, not external-library comparisons. One NTT natural-input/bit-reversed-output cell is unavailable and is omitted rather than encoded as zero.}
\end{table*}
